% !TeX root = ../thuthesis-example.tex

\begin{acknowledgements}
  行文至此，虽然由于种种原因导致三个月后的去向依然扑朔迷离，但是仍然想要好好感谢一路以来遇到的良师益友。
  
  首先要感谢我的导师赵慧婵老师。在大三下学期与赵老师相识，当时刚从加拿大交换回来，一切又都是重新的开始，是赵老师把我带进实验室开始做一些工作。感谢从那时起，赵老师一直以来对我的关心和指导，无论是科研、生活，乃至对于未来去向的建议和指导。

  感谢我父母一直以来的关心和支持。没有你们的帮助，我没有办法如此顺利的进入清华大学学习，结实这么多良师益友，为自己的未来打下基础。感谢从小时候起无微不至的关心与照顾，并且一路上给予我鼓励、支持和宝贵的建议，让我能够更加坚定地走下去！

  还要感谢课题组的师兄师姐们。感谢朱颖雷师兄邀请我一起合作某一个项目，同时对我工作提供仿真方面的建议和帮助；感谢汪田鸿博士、周亮博士、陈信先师兄在器件制造方面提供的指导与支持；感谢王启尧师兄、杜博源博士、於思晓同学在我刚加入课题组时帮助我融入；感谢潘依依博士、邵琦博士对项目整体结构的指导和帮助。

  感谢所有帮助过我的老师们。感谢吴丹老师方方面面的支持和帮助，从交换期间的资金支持，到大三暑期的实习，申请期间的大力推荐；感谢汪泽老师在我初入清华时，对我进行科研的训练和指导；感谢尹文生老师在我申请期间的支持和帮助。感谢多伦多大学的 Hugh H. T. Liu（刘泓涛）院士，是您在我初到加拿大时接纳了我，并且把我带向无人机相关领域，也让我了解到了国内外科研的异同。

  还有很多很多要感谢的同学。感谢清华交响的王天天、王圳豪、任益学、邹箫桐等等，是你们的出现让我的大学生活更加丰富多彩；感谢机械系的郭任杰、张峻瑜、陈诺博士等等，有你们在我才更好的加入了机械系这个大家庭；感谢UTIAS的秦超博士、钱龙昊博士对我在FSC课题组时的指导和帮助。感谢小米科技有限公司接纳我实习，感谢殷傲同事、苏炜同事、许中科经理对我实习期间无论是生产还是生活方面的关心。

  感谢江亚羲同学。每当我遇到挫折的时候都会坚定的支持我，从暑研签证受挫，到暑期实习更换单位、申请时的迷茫、毕设快要结项时遭遇的一系列“黑天鹅”事件，都与我一起思考解决方案，帮我克服困难。

  最后，感谢参与论文评阅的老师们在百忙之中为我提出宝贵意见。
\end{acknowledgements}
